\documentclass[12pt, french, latin.medieval, a4paper]{report}
% package de langues 
\usepackage[french]{babel}
\usepackage[utf8]{inputenc}
\usepackage[T1]{fontenc}

% package pour référence de liens sur pdf 
\usepackage[pdftex]{hyperref} % liens
\usepackage{lastpage} % numéro de page
\usepackage{caption}

% package de formes de document 
\usepackage[left=2.5cm,
              right=2.5cm,
              top=2.5cm,
              bottom=2.5cm
             ]{geometry}
% pour les tableaux 
\usepackage{array,multirow,makecell} 
\newcolumntype{R}[1]{>{\raggedleft\arraybackslash }b{#1}}
\newcolumntype{L}[1]{>{\raggedright\arraybackslash }b{#1}}
\newcolumntype{C}[1]{>{\centering\arraybackslash }b{#1}}

% package pour la page de couverture 
\usepackage{graphicx}
\usepackage[cyr]{aeguill}
\usepackage{fancyheadings}
\usepackage{amsmath}
\usepackage{amsfonts}
\usepackage{amssymb}
\usepackage{geometry}
\usepackage{eso-pic}
\usepackage{tikz}
\usetikzlibrary{calc}


% nouvelle commande pour la forme 
%\fancyhf[OFC]{\thepage / \pageref{LastPage}}
\rfoot{\thepage / \pageref{LastPage}}
%\setcounter{tocdepth}{3}

% pakages de couleurs du document
\usepackage{color}
\usepackage{xcolor}


% package pour la chimie
\usepackage[version=4]{mhchem}


% package pour la visualisation du code
\usepackage{listings}
\definecolor{darkWhite}{rgb}{0.94,0.94,0.94}
 
\lstset{
  aboveskip=3mm,
  belowskip=-2mm,
  backgroundcolor=\color{darkWhite},
  basicstyle=\footnotesize,
  breakatwhitespace=false,
  breaklines=true,
  captionpos=b,
  commentstyle=\color{red},
  deletekeywords={...},
  escapeinside={\%*}{*)},
  extendedchars=true,
  framexleftmargin=16pt,
  framextopmargin=3pt,
  framexbottommargin=6pt,
  frame=tb,
  keepspaces=true,
  keywordstyle=\color{blue},
  language=python,
  literate=
  {²}{{\textsuperscript{2}}}1
  {⁴}{{\textsuperscript{4}}}1
  {⁶}{{\textsuperscript{6}}}1
  {⁸}{{\textsuperscript{8}}}1
  {€}{{\euro{}}}1
  {é}{{\'e}}1
  {è}{{\`{e}}}1
  {ê}{{\^{e}}}1
  {ë}{{\¨{e}}}1
  {É}{{\'{E}}}1
  {Ê}{{\^{E}}}1
  {û}{{\^{u}}}1
  {ù}{{\`{u}}}1
  {â}{{\^{a}}}1
  {à}{{\`{a}}}1
  {á}{{\'{a}}}1
  {ã}{{\~{a}}}1
  {Á}{{\'{A}}}1
  {Â}{{\^{A}}}1
  {Ã}{{\~{A}}}1
  {ç}{{\c{c}}}1
  {Ç}{{\c{C}}}1
  {õ}{{\~{o}}}1
  {ó}{{\'{o}}}1
  {ô}{{\^{o}}}1
  {Õ}{{\~{O}}}1
  {Ó}{{\'{O}}}1
  {Ô}{{\^{O}}}1
  {î}{{\^{i}}}1
  {Î}{{\^{I}}}1
  {í}{{\'{i}}}1
  {Í}{{\~{Í}}}1,
  morekeywords={*,...},
  numbers=left,
  numbersep=10pt,
  numberstyle=\tiny\color{black},
  rulecolor=\color{black},
  showspaces=false,
  showstringspaces=false,
  showtabs=false,
  stepnumber=1,
  stringstyle=\color{gray},
  tabsize=4,
  title=\lstname,
}

%package pour les algorithmes
\usepackage{algorithm}
\usepackage{algorithmic}
\usepackage{amsmath}

% package pour les tableaux
\usepackage{longtable}
\usepackage{booktabs}

% package pour les figures 
\usepackage{subcaption}

% package de reférencement
% \usepackage[backend=biber,safeinputenc, sorting=none]{biblatex}

% mettre la bibliographie dans la table de matières
%\usepackage[nottoc, notlof, notlot]{tocbibind}
\include{page_couverture.tex}


\author{MBE MBE MINDJANA LOIC HENRI}
\title{GNPS et MDByTE}

\begin{document}
    
    \pagecouverture
    {Representation Learning in Multi-view Clustering : A Literature Review}
    {Man-Sheng Chen et al.}
    \tableofcontents
    \vspace{20cm}
    
    \section{Introduction}
    De plus en plus de nos jours les données sont obtenues de diverses sources qui sont généralement liées à des tâches différentes sur des objets communs. Par exemple, les données sur les voitures qu'on retrouvent généralement sur des sites sont leurs images et descriptions textuelles qui constituent des représentations différentes (\textbf{des vues}) de ces dernières. Afin de mieux explorer ce type de données, on utilise des algorithmes de regroupement spécifiques appelés \textbf{Algorithmes de Clustering Multi-vues} qui permettent de tenir compte de la complémentarité et du consensus (accord) entre les différentes perspectives des données communément appelées \textbf{vues}. Comme les chercheurs \textbf{Yan Yang} et \textbf{Hao Wang}, Man-Sheng Chen, Jia-Qi Lin et al. ont fait une enquête visant à recenser l'ensemble de ces algorithmes et à définir une taxonomie. Ce recensement a conduit à l'article de référence \cite{chen2022representation} qui sera le sujet de ce document. Dans les lignes qui suivent nous présenterons d'abord l'article ensuite nous parlerons la taxonomie des algorithmes de clustering Multi-vues tel que décrit dans l'article et enfin nous listerons des datasets y présentées et les orientations de recherche brûlantes selon les auteurs.
    
    \section{Description de l'article}
	L'article \textbf{Representation Learning in Multi‑view Clustering: A Literature Review} \cite{chen2022representation} est le résultat d'une enquête menée plus récente que celle de l'article \cite{yang2018multi}, sur les algorithmes de clustering multi-vues fait par les chercheurs Man-Sheng Chen, Jia-Qi Lin et al. Dans cet article, les auteurs présentent une typologie des algorithmes de clustering multi-vues basée sur l'apprentissage de représentation et les récentes directions de recherche sur le sujet du clustering multi-vues. on as les informations suivantes sur l’article:
	\begin{itemize}
    	\item \textbf{Titre}: Representation learning in multi-view clustering: A literature review
    	\item \textbf{Auteurs}:  Man-Sheng Chen, Jia-Qi Lin , Xiang-Long Li, Bao-Yu Liu, Chang-Dong Wang, Jian-Huang Lai, Dong Huang
    	\item \textbf{Année de publication}: 2022
    	\item \textbf{Journal}: Data Science and Engineering
	\end{itemize}

	\section{Clustering multi-vues}
	Le clustering multi-vues est une technique d'apprentissage non supervisée qui consiste à regrouper des données(observations) similaires entre eux et d'en éloigner ceux qui ne le sont pas en tenant compte qu'elles sont constituées de différentes vues. Les vues sont perçues comme les différentes représentations des observations utilisées.

	On peut séparer actuellement les algorithmes de clustering multi-vues en 2 catégories tel que sur la figure \ref{fig:typo-mvc}, ceux qui apprennent une représentation unifiée des différentes vues des données ( on parle de \textbf{"Representation Learning based MVC"} en anglais) et ceux qui ne le font pas (on parle de \textbf{"Non-representation Learning based MVC"} en anglais). Nous présenterons dans les prochaines sections les différents types d'algorithmes pour ces catégories.

	\begin{figure}
    	\centering
    	\includegraphics[width= \linewidth]{Images/taxonomyMultiViewClustering.png}
    	\caption{Taxonomie des algorithmes de clustering multi-vues \cite{chen2022representation}}
    	\label{fig:typo-mvc}
	\end{figure}

    \section{Algorithmes sans apprentissage de représentation}
        On peut citer :
        \subsection{Algorithmes de co-formation}
         	Dans cette catégorie d'algorithme, on cherche à maximiser l'accord entre les vues tout en améliorant les clusters itérativement de vues en vues. On construit initialement les clusters pour chaque vue, et les clusters d'une vue sont améliorés en utilisant ceux d'autres vues dans le but de maximiser le consensus entre eux. On procède ainsi itérativement tour à tour pour chaque vue jusqu'à obtenir une convergence du consensus entre les vues ou après un certain nombre d'itérations.
       	 
        	La figure \ref{fig:co-train1} l'illustre assez bien. On construit des clusters pour chacune des vues 1 et 2 après entraînement à la suite on améliore les clusters d'une vue, en utilisant des connaissances de l'autre dans le but maximiser la cohésion entre les vues.
       	 
        	\begin{figure}
            	\centering
            	\includegraphics[width=\linewidth]{Images/co-training.png}
            	\caption{Processus général des algorithmes de co-formation \cite{yang2018multi}}
            	\label{fig:co-train1}
        	\end{figure}
        	On peut dès lors comprendre que ces algorithmes possèdent un défaut majeur qui est le temps d'exécution, car on reconstruit les clusters en un certain nombre d'itérations.
    
        \subsection{Algorithmes multi-noyaux}
        	S'il est difficile de construire des clusters dans un espace d'une dimension donnée, on pourrait plus aisément le faire en projetant les observations dans un espace correspondant de plus grande dimension. C'est sur ce principe que les algorithmes de clustering par noyaux sont construits pour les données représentées par des caractéristiques d'une même représentation. On recherche donc une fonction noyau qui permettra d'obtenir de meilleurs clusters suivant un critère d'évaluation comme l'inertie intra-classe.
       	 
        	Pour cette catégorie de clustering multi-vues, on cherche également une fonction noyau qui permettra d'obtenir des clusters dans un espace de plus grande dimension qui agglomère les espaces des différentes vues. Ceci est fait en combinant linéairement ou non linéairement des fonctions noyaux obtenues, sur chacunes des vues dans le but d'obtenir des clusters en tenant compte le plus possible des principes de complémentarité et de consensus donnés plus haut. La figure \ref{fig:multi-kernel1} décrit assez bien le processus global des algorithmes de ce type, on a des vues $View_1,View_2,...,View_m$, dont on obtiendra des fonctions noyaux correspondantes qui seront combinées linéairement ou non linéairement par apprentissage pour obtenir un noyau unifié.
    
        	Généralement ces algorithmes s'inspirent d'autres algorithmes de clustering à une vue tel que KMEANS, et les modifient pour intégrer les fonctions noyaux dans le contexte des données multivues. Dans le cas de KMEANS, on définit la fonction noyau unifiée $K$ comme dans l'équation \ref{equa:kernel1}. vu qu'on souhaite minimiser l'inertie intraclasse du clustering $I$, on aura le problème d'optimisation \ref{equa:kernel2} où $\phi$ est la fonction de projection dans une grande dimension et $W$ est la matrice de poids permettant d'exploiter les deux principes de consensus et de complémentarité. Il faut noter que $\gamma$ est la matrice indicateur de clusters c'est à dire que si l'observation $i$ est dans le cluster $k$ alors $\gamma_{ik}=1$ sinon $\gamma_{ik}=0$.
                \begin{equation}
                    K(x,y) = {[\exp(-\frac{{||x-y||}^2}{2 \sigma^2}) + R]}^P \\
                    \text{ avec $R$ et $P$ sont des réeels }
                    \label{equa:kernel1}
                \end{equation}
        
                \begin{equation}
                    \begin{aligned}
                        \min_{W,R,P} I = \sum_{v=1}^{V} \sum_{i=1}^{n} w_{i,v} \sum_{k=1}^{K} \gamma_{ik} {|| \phi(x_{i}^{v}) - \phi(c_{j}^{v}) ||}^2 
                        \\
                        = \sum_{v=1}^{V} \sum_{i=1}^{n} w_{i,v} \sum_{k=1}^{K} 2 \gamma_{ik} [{(1 + R)}^P + K(x_{i}^{v}, c_{j}^{v})]
                        \\ 
                        \text{s.t, } w_{i,v} \geq 0\text{, }\prod_{v} w_{i,v}=1
                    \end{aligned}
                    \label{equa:kernel2}
                \end{equation}
                
                \begin{figure}
                    \centering
                    \includegraphics[width = \linewidth]{Images/multi-kernel.png}
                    \caption{Processus général des algorithmes multi-noyaux \cite{yang2018multi}}
                    \label{fig:multi-kernel1}
                \end{figure}
                
        \subsection{Algorithmes de factorisation matricielle}
        	A la base étant une technique de réduction de dimensionnalité, la factorisation de données est grandement utilisée pour résoudre des problèmes de recommandation et de classification. Ils sont utilisés pour le problème de clustering multi-vues en factorisant les matrices de données des vues (non négatives) en des couples de matrices dit facteurs. Pour chacun de ces couples, on a l'un qui représente l'espace latent lié à une vue précise et l'autre l'indicateur des clusters pour cette même vue.
       	 
        	L'indicateur unifié des clusters est ainsi calculé par agglomération des indicateurs des couples de facteurs obtenus en minimisant l'erreur de factorisation et d'agglomération des indicateurs de chaque vue. Plus formellement, on a généralement le problème d'optimisation de l'équation \ref{equa:nmf}, où $U^{(v)}$ désigne le facteur latent de la vue $v$, $C^{(v)}$ le facteur indicatif de cluster de la vue $v$, $C^*$ est l'indicateur unifié ou dit de consensus globale des vues et $\Phi$ est une fonction de fusion de $C^*$ et $C^{(v)}$ qui traduit une erreur d'agglomération. Il faut noter que la norme $||A||_{F}^{2} = tr(A.A_{T}) = \sum_{i,j} a_{i,j}^2$ est appelée la norme de Frobenius

        	\begin{equation}
             	\begin{aligned}
                	\min {\sum_{v=1}^{V} {|| X^{(v)} - U^{(v)}.C^{(v)} ||}_F^2 + \Phi(C^{(v)}, C^{*})} \\
                	\text{s.t, $U^{(v)},C^{(v)},C^{*} \geq 0$,}
            	\end{aligned}
            	\label{equa:nmf}
        	\end{equation}
            
        \subsection{Algorithmes spectraux}
            Effectuer du clustering revient à trouver k groupes d'observations similaires entre eux. Cela se fait par une fonction $f$ qui permet d'attribuer à chaque observation un indicateur propre à un des $K$ clusters. On peut donc voir que $f$ est lui même un indicateur pour attribuer une observation à un cluster donné et donc peut être vu comme une matrice $(n \times k)$ où $f_{i,k}$ est la probabilité que l'observation $i$ soit dans la classe $k$.
            
            L'objectif des algorithmes de clustering de ce type est de calculer cette matrice tout en tenant compte des similarités entre les observations(par ricochet d'un graphe de similarité), en cherchant à résoudre le problème d'optimisation \ref{equa:spectral1} où $L$ est le laplacien du graphe de similarité des observations.
        
            \begin{equation}
                \min_{F^{T}F = I} tr(F^{T}LF)
                \label{equa:spectral1}
            \end{equation}


            Pour pouvoir prendre en compte plusieurs vues, le problème d'optimisation \ref{equa:spectral1} change pour le problème \ref{equa:spectral2} qui est équivalent au problème \ref{equa:spectral3} où $L = \sum_{v=1}^{V} \alpha_{v}^{r} \times L_{v}$, avec $r$ étant un coefficient pour contrôler les distributions des vues. Nous remarquons donc que ce problème est à peu près identique au problème de clustering spectral à une vue.
    
            \begin{equation}
                \begin{aligned}
                    \min_{F^{T}F = I, \alpha} \sum_{v=1}^{V} \alpha_{v}^{r} tr(F^{T}L_{v}F) \\
                    \text{s.t, } \sum_{v} \alpha_{v} = 1 \text{ , } \alpha_{v} \geq 0
                \end{aligned}
                \label{equa:spectral2}
            \end{equation}
    
            \begin{equation}
                \begin{aligned}
                    \min_{F^{T}F = I, \alpha} tr(F^{T}LF) \\
                    \text{s.t, } \sum_{v} \alpha_{v} = 1 \text{ , } \alpha_{v} \geq 0
                \end{aligned}
                \label{equa:spectral3}
            \end{equation}
    
            Les vecteurs propres dans l'ordre croissant de leurs valeurs propres du laplacien conjugué $L$ seront calculés, puis soumis à un algorithme de clustering classique pour pouvoir trouver $F$. Il faut noter qu'avec cette méthode de clustering, il est possible de détecter le nombre de clusters approprié (qui est le nombre de vecteurs propres), en s'appuyant sur les écarts entre les valeurs propres consécutives. Le nombre $K$ de clusters peut être choisi ou obtenu en détectant une grande variation entre deux valeurs propres consécutives selon l'ordre croissant de ces valeurs.

	\section{Algorithmes basés sur l'apprentissage de représentation}
    	Dans cette catégorie d'algorithmes, on a ceux qui utilisent les réseaux de neurones pour construire les représentations (on parle de "Deep representation learning-based MVC" en anglais) et d'autres qui ne les utilisent pas (on parle de "Shallow representation learning-based MVC" en anglais) dit d'algorithmes peu profond.

    	\subsection{Algorithmes basés sur l'apprentissage peu profond de représentation }
        	\subsubsection{Algorithmes basés sur les graphes}
        	Pour un ensemble d'observations n'ayant qu'une vue, on peut construire un graphe de similarité(très proche d'un réseau moléculaire) qui permettrait après découpe (clustering spectral) d'obtenir des clusters.
       	 
        	L'idée \ref{fig:graph1} pour pouvoir tenir compte des données à plusieurs vues, est de construire pour chaque vue des graphes de similarités. Puis de définir une politique de fusion de ces graphes, qui s'appuie généralement sur l'attribution des poids à des observations par rapport aux groupes formés. De là on obtient un graphe unifié de similarité qui pourra être sujet d'un clustering spectral pour pouvoir construire des clusters. On cherchera ainsi à trouver les poids amenant à avoir le moins d'erreur possible de fusion des vues.
    
        	\begin{figure}
            	\centering
            	\includegraphics[width=\linewidth]{Images/graph2.png}
            	\caption{Processus général des algorithmes basés sur les graphes \cite{chen2022representation}}
            	\label{fig:graph1}
        	\end{figure}

        	Plus formellement le problème d'optimisation générale est celui de l'équation \ref{equa:graph1}, où $A^{(v)}$ est la matrice d'affinité spécifique à la vue $v$, $S$ représente la matrice de similarité du graphe unifié, $Dis$ est utilisé pour représenter certaines métriques de similarité ou de dissimilarité utilisées dans les vues et $\Phi$ est une fonction d'erreur de fusion de $S$ et $A^{(v)}$.

        	\begin{equation}
             	\begin{aligned}
                	\min {\sum_{v=1}^{V} \sum_{i \neq j}^{n} \text{Dis}(X_i^{(v)}, X_j^{(v)})A^{(v)} + \Phi(C^{(v)}, C^{*})} \\
                	\text{s.t, $U^{(v)},C^{(v)},C^{*} \geq 0$,}
            	\end{aligned}
            	\label{equa:graph1}
        	\end{equation}
            
        	\subsubsection{Algorithmes de sous-espaces}
       	 
        	L'idée ici est de rechercher une représentation unifiée des vues pour ensuite y effectuer un algorithme de clustering à une vue. De ce fait on procède généralement de deux manières tel que sur la figure \ref{fig:subspace}:
        	\begin{enumerate}
            	\item soit on apprend directement la représentation depuis les sous-espaces des vues des données.
            	\item soit on apprend d'abord un espace latent avant de l'utiliser pour apprendre la représentation unifiée. Cela est fait sur la base des sous espaces des vues   
        	\end{enumerate}

        	Globalement, pour chaque vue $v$, on cherche par résolution du problème d'optimisation \ref{equa:subspace1}, les matrices $Z^{(v)}$ traduisant les données dans le sous-espace de la vue $v$ et $Z^*$ la représentation commune unifiée des observations des vues obtenues par les matrices de sous-espaces $Z(v)$. On as dans cette équation $\Omega(.)$ étant une fonction de régularisation de l'apprentissage de $Z^{(v)}$ et $\phi$ une fonction d'erreur de fusion de $S$ et de $Z^{(v)}$. Il faut noter que $\lambda$ et $\gamma$ sont des termes constants à définir, qui permettent de quantifier à quel point on tiendra compte de la régularisation et de la fusion respectivement.
    
        	\begin{figure}
            	\centering
            	\includegraphics[width = \linewidth]{Images/subspace.png}
            	\caption{Processus général des algorithmes de sous-espaces \cite{yang2018multi}}
            	\label{fig:subspace}
        	\end{figure}
    
        	\begin{equation}
            	\begin{aligned}
                	\min {\sum_{v=1}^{V} {||X^{(v)} - X^{(v)}Z^{(v)}||}_F^2 + \lambda\Omega(Z^{(v)}) + \gamma\Phi(Z^{*},Z^{(v)}) \\
                	\text{s.t,$Z^{(v)} \geq 0$, $Z^{(v)}\textbf{1}=1$}
            	\end{aligned}
            	\label{equa:subspace1}
        	\end{equation}

        	Pour apprendre un espace latent $H$, on résout par contre le problème d'optimisation \ref{equa:subspace2} où $\Theta$ est une fonction évaluant l'erreur d'utilisation des données de la vue $v$ pour obtenir la représentation unifié $Z^*$. Dans cette équation $\gamma$ est une constante permettant de réguler l'erreur du calcul de l'espace latent sur la base de $Z^*$.
       	 
        	\begin{equation}
            	\begin{aligned}
                	\min \sum_{v=1}^{V}\Theta(X^{(V)},Z^*) + \gamma{||H - HZ^*||}_F^2 \\
                	\text{s.t,$Z^* \geq 0$, $Z^*.\textbf{1}=1$}
            	\end{aligned}
            	\label{equa:subspace2}
        	\end{equation}
            
    	\subsection{Algorithmes basés sur l'apprentissage profond de représentation}
        	Dans le but d'explorer complètement les espaces de données originaux pour construire une représentation unifiée, on pourrait utiliser des réseaux de neurones au vue de leur performance sur des problèmes supervisés. De cette idée les algorithmes de clustering multi-vues utilisant les modèles de deepLearning ont été créés. Comme le modèle d'auto-encodage \cite{wang2016auto} utiliser dans le domaine du traitement du langage naturel, ces réseaux sont entraînés par rétropropagation des erreurs liées aux regroupements effectués sur les représentations obtenues des observations. La figure \ref{fig:deepLearning} illustre très bien le processus général de ces algorithmes.

        	\begin{figure}
            	\centering
            	\includegraphics[width = \linewidth]{Images/multiviewCLusteringDeep.png}
            	\caption{Processus général des algorithmes d'apprentissage profond de représentation \cite{chen2022representation}}
            	\label{fig:deepLearning}
        	\end{figure}
	\section{Datasets recensés}
	Les auteurs au cours de leur enquête, ont recensé un ensemble de bases de données d'entraînement multi-vues, qui nous ont interpellés. Elles sont:
	\begin{itemize}
    	\item La dataset ACM contenant un graphe de 3025 articles, malheureusement aucune référence a été donné
    	\item La dataset \href{http://mlg.ucd.ie/datasets/segment.html}{BBC4view} contenant des actualités présentés par la BBC.
    	\item La dataset \href{http://cvc.yale.edu/projects/yalefaces/yalefaces.html}{Yale}  contenant des images de visages avec différents types de caractéristiques extraites de ces dernières
	\end{itemize}

	\section{Orientations de recherche}
	Les auteurs ont recensé les différentes orientations de recherche pour le futur, compte tenu des différentes perspectives et résultats des articles qu'ils ont lus, et d'un certain nombre de problèmes qui ont été levés. on peut citer comme domaine:
	\begin{itemize}
    	\item \textbf{Le clustering multi-vues à grande échelle} : Sachant que dans les applications du monde réel, il est très courant qu'on manipule des données de l'ordre du millions. Il est donc important que les algorithmes de clustering multi-vues puissent passer à l'échelle. Des tentatives pour résoudre ce problème ont été proposées, principalement dans la branche des algorithmes d’apprentissage de représentations. Nous pensons que la parallélisation des algorithmes de clustering multi-vues pourrait aider à résoudre le problème de passage à l’échelle.
    	\item \textbf{Le clustering multi-vues basé sur l'apprentissage profond} : Généralement on utilise les modèles d'apprentissage profond pour des problèmes d'apprentissage supervisés du fait de leur performance. Des recherches ont été faites pour les algorithmes dit de représentation pour les introduire dans le processus de calcul de représentation. Malheureusement selon les auteurs les algorithmes résultant de ces recherches ne permettent pas encore de bien représenter l’espace unifié des vues.

    	\item \textbf{Le clustering multi-vues de données incomplètes} : On obtient dans la vie courante les données sur des objets au fur et à mesure, qu’on en trouve. Ce qui nous amène généralement dans les cas où pour un objet, on n'a pas toutes les données liées à toutes ses vues pertinentes. comment effectuer du clustering avec de telles données est une question très pertinente. Avoir des algorithmes gérant de telles données permettrait de réaliser des regroupements de plus en plus précis au fur et à mesure qu'on recueille des données sur les vues des observations.
	\end{itemize}

	\section{Conclusion}
	En somme, il était question pour nous de présenter la taxonomie des algorithmes de clustering multi-vues donnée par les auteurs, présenter les datasets données qui nous ont le plus parlé et parler des orientations de recherche brûlant dans ce domaine. Compte tenu des données spectrales qu'on possède, nous pensons qu'on pourrait utiliser l'apprentissage de représentation unifié, pour exploiter la complémentarité des spectres MS et RMN. Nous pensons qu'utiliser l'algorithme de clustering multi-vues de graphe est idéale compte tenu qu'on as en entrée des graphes des observations et qu'on recherche une représentation de consensus et de complémentarité sans pour autant impliquer les regroupements des observations.

   
\bibliographystyle{plain}
    \bibliography{reference_GNPS_MADByTE}
\end{document}